% =====================================================================
%  cv-mphil-tmr.tex  --  TAILORED CV
%  Cambridge MPhil in Translational Medical Research (Dept of Medicine),
%  Michaelmas 2027 entry, regular round (no Gates Cambridge application).
%
%  BEFORE SUBMITTING: Cambridge's postgraduate admissions rules say "AI
%  must not be used when writing personal statements \& CVs", including
%  help with English. The entries below were drafted with Claude. Rewrite
%  every entry in your own words before this goes into the portal; the
%  layout (cvstyle.sty) is typesetting \& can stay.
%  https://www.postgraduate.study.cam.ac.uk/apply/how/ai-postgraduate-applications
%
%  WHAT THE COURSE WILL CHECK, \& WHERE THIS CV ANSWERS IT
%  (Cambridge Careers' three tests: eligibility, suitability, potential)
%
%  ELIGIBILITY -> Education, first
%   * "Good II.i" required. Cambridge's US table: 3.5 = 2:1,
%     3.7 = High 2:1, 3.9 = First. The transcript GPA is 3.60, which
%     meets the requirement as a 2:1, \& the CV says so ("UK 2:1
%     equivalent", set in cvdata.tex).
%   * Modules are a short block, not bullets, with NO grades: the science
%     grades are mixed \& the transcript is uploaded anyway. The block
%     lists every related module, alphabetically, to show coverage
%     ("significant relevant preparation").
%   * Bold (\key) rule: a module is bold only if it is relevant to the
%     course \& has the top grade (NJIT A; Lund 5 or VG), because a
%     reader who sees bold will look that module up on the transcript.
%     No Lund module has Lund's top grade, so none is bold; the Lund
%     modules are split by level instead (four are master's level).
%     Bold now: Biomedical Ethics, Medicine \& Law in Modern History
%     (research ethics, regulation, patient involvement); Linear Algebra,
%     Mathematical Modeling (data analysis \& modelling core); the
%     amyloid modelling independent study (MATH 491, A). Not bold despite
%     relevance: Cell \& Molecular Biology (B), Biochemistry (C),
%     Probability \& Statistics (B+).
%   * Independent study row: research done as a self-contained
%     one-semester course is listed here, not under Research Experience.
%     Credit earned inside a longer project stays under Research:
%     BME 491 (Summer 2025) was KumarLab work; MATH 310 Co-op (Summer
%     2026) was the Imperial placement; BMET 311-312 (Applications in
%     Healthcare I-II) was the EMS study; HIST 403 is credit for the
%     Harvard Sikh Center fellowship, which is listed under Funding.
%   * Merit scholarships \& academic prizes go inside the Education
%     entry (as in the Cambridge sample CVs), not in a separate list.
%
%  POTENTIAL -> Research, then Outputs, then Funding
%   * Half the course is a 32-week project in a host lab. It is assessed
%     by a 5,000-word literature review, which may include a section
%     written like a GRANT PROPOSAL, \& a 5,000-word results write-up,
%     plus an oral presentation (Lent) \& a poster (Easter). The
%     P\&emic Preparedness theme runs a mock GBP 2m funding pitch.
%   * So: research first (can you run a project), publications \&
%     presentations next (can you write it up \& present it), then
%     funding (others have already judged your proposals worth funding).
%
%  SUITABILITY -> Clinical \& Translational Experience
%   * Course aims name: clinical trials, bench-to-bedside, the patient-
%     physician interface, the academic-pharma interface, NHS research,
%     \& patient & public involvement (PPIEP workshop). Themes:
%     translational immunology, p\&emic preparedness, cardiovascular &
%     respiratory, acute & emergency care (workshop: "How to run a
%     clinical study"). The heading uses the course's own words.
%   * Cambridge Careers: give a key requirement its own heading (their
%     sample separates "Commercial experience" for that reason).
%
%  BREADTH -> Leadership, Teaching \& Mentoring
%   * Not scored for itself in the regular round. It shows communication
%     \& mentoring, \& gives interviewers something to ask about. It is
%     the first section to shrink if space is needed.
%
%  LOOK-UP BLOCK -> Research Skills \& Training, last
%   * Skills are claims; the research bullets are the evidence, so the
%     evidence comes first \& this block works as an index. It still
%     matters: students rank five host groups in the summer, \& the
%     FAQ says lab time depends on whether a project is wet or dry.
%     Past hosts include dry/hybrid groups (ML drug discovery, cell-atlas
%     data science, AI in dementia, perioperative data science), so name
%     computational tools precisely.
%
%  WHO READS IT
%   Course Director Dr Harriet Groom (CITIID); Academic Director Prof.
%   Yorgo Modis (structural biology); theme leads Cook, R\&ow,
%   Matheson, Rihn, Jorgensen, Clarke, Newcombe, Conway Morris.
%   Interviews: 2-3 of the course team, ~20 min, opening with a 3-minute
%   talk on motivation \& relevant experience, then questions \& a
%   critical-thinking question. Anything here can be asked about.
%
%  LEFT OUT ON PURPOSE
%   * Personal profile: the three text boxes (Reasons 1000 characters,
%     Relevant Experience 1000, Career Goals 1500) do that job, \&
%     Cambridge Careers calls profiles "generally unnecessary" in the UK.
%   * "Additional"/"Other" experience, soft-skills lists, bare
%     memberships: they do not evidence anything the course checks.
%   * References: entered in the portal (2 referees).
%
%  RULES: two pages max; most relevant item first within a section; every
%  bullet = what you did + how + result, with a number where possible;
%  no section split across the page break.
%
%  IF OVER TWO PAGES, CUT IN THIS ORDER
%   Industry sub-group (fold into a research bullet) -> presentations to
%   2 -> leadership to 3 -> independent study row to the completed item.
%   Global Studies track \& the full exchange entry are already cut
%   (both are on the general CV).
% =====================================================================
\documentclass[a4paper,11pt]{article}  % drop to 10pt only if needed to hold two pages; never below 10pt
\usepackage{cvstyle}
% =====================================================================
%  cvdata.tex  --  facts shared by every version of the CV
%  Each version loads this with % =====================================================================
%  cvdata.tex  --  facts shared by every version of the CV
%  Each version loads this with % =====================================================================
%  cvdata.tex  --  facts shared by every version of the CV
%  Each version loads this with \input{cvdata} after \usepackage{cvstyle}.
%  Change a fact here once and every version updates.
% =====================================================================
\renewcommand{\cvname}{Kabir Singh}
\renewcommand{\cvupdated}{October 2026}

\newcommand{\cvemail}{\href{mailto:kabirsinghji@gmail.com}{kabirsinghji@gmail.com}}
\newcommand{\cvphone}{+1 917 636 8664}
\newcommand{\cvweb}{\href{https://kabirsinghji.github.io}{kabirsinghji.github.io}}
\newcommand{\cvorcid}{\href{https://orcid.org/0009-0002-5159-5879}{0009-0002-5159-5879}}
\newcommand{\cvscholar}{\href{https://scholar.google.com/citations?user=-RpJmPsAAAAJ}{Google Scholar}}
\newcommand{\cvlinkedin}{\href{https://www.linkedin.com/in/kabir-singh-ji}{linkedin.com/in/kabir-singh-ji}}

% ---- GPA --------------------------------------------------------------
% \gpaNow  NJIT unofficial transcript of 14 Oct 2026 (includes Summer
%          2026, HIST 213 = A): 429.50 quality points / 119 graded credits
%          = 3.609, printed 3.61. MATH 310 Co-op Work Experience I is still
%          in progress and Fall 2026 is not graded yet.
% \gpaUK   Used only in the UK-facing versions. Cambridge's published
%          table for US degrees: 3.5 = 2:1, 3.7 = High 2:1, 3.9 = First.
%          3.60 is therefore a 2:1. Change this only if the GPA crosses 3.7.
%          To drop the label, delete "(\gpaUK)" where it is used.
\newcommand{\gpaNow}{3.61}
\newcommand{\gpaUK}{UK Good 2:1 Honours equivalent}
 after \usepackage{cvstyle}.
%  Change a fact here once and every version updates.
% =====================================================================
\renewcommand{\cvname}{Kabir Singh}
\renewcommand{\cvupdated}{October 2026}

\newcommand{\cvemail}{\href{mailto:kabirsinghji@gmail.com}{kabirsinghji@gmail.com}}
\newcommand{\cvphone}{+1 917 636 8664}
\newcommand{\cvweb}{\href{https://kabirsinghji.github.io}{kabirsinghji.github.io}}
\newcommand{\cvorcid}{\href{https://orcid.org/0009-0002-5159-5879}{0009-0002-5159-5879}}
\newcommand{\cvscholar}{\href{https://scholar.google.com/citations?user=-RpJmPsAAAAJ}{Google Scholar}}
\newcommand{\cvlinkedin}{\href{https://www.linkedin.com/in/kabir-singh-ji}{linkedin.com/in/kabir-singh-ji}}

% ---- GPA --------------------------------------------------------------
% \gpaNow  NJIT unofficial transcript of 14 Oct 2026 (includes Summer
%          2026, HIST 213 = A): 429.50 quality points / 119 graded credits
%          = 3.609, printed 3.61. MATH 310 Co-op Work Experience I is still
%          in progress and Fall 2026 is not graded yet.
% \gpaUK   Used only in the UK-facing versions. Cambridge's published
%          table for US degrees: 3.5 = 2:1, 3.7 = High 2:1, 3.9 = First.
%          3.60 is therefore a 2:1. Change this only if the GPA crosses 3.7.
%          To drop the label, delete "(\gpaUK)" where it is used.
\newcommand{\gpaNow}{3.61}
\newcommand{\gpaUK}{UK Good 2:1 Honours equivalent}
 after \usepackage{cvstyle}.
%  Change a fact here once and every version updates.
% =====================================================================
\renewcommand{\cvname}{Kabir Singh}
\renewcommand{\cvupdated}{October 2026}

\newcommand{\cvemail}{\href{mailto:kabirsinghji@gmail.com}{kabirsinghji@gmail.com}}
\newcommand{\cvphone}{+1 917 636 8664}
\newcommand{\cvweb}{\href{https://kabirsinghji.github.io}{kabirsinghji.github.io}}
\newcommand{\cvorcid}{\href{https://orcid.org/0009-0002-5159-5879}{0009-0002-5159-5879}}
\newcommand{\cvscholar}{\href{https://scholar.google.com/citations?user=-RpJmPsAAAAJ}{Google Scholar}}
\newcommand{\cvlinkedin}{\href{https://www.linkedin.com/in/kabir-singh-ji}{linkedin.com/in/kabir-singh-ji}}

% ---- GPA --------------------------------------------------------------
% \gpaNow  NJIT unofficial transcript of 14 Oct 2026 (includes Summer
%          2026, HIST 213 = A): 429.50 quality points / 119 graded credits
%          = 3.609, printed 3.61. MATH 310 Co-op Work Experience I is still
%          in progress and Fall 2026 is not graded yet.
% \gpaUK   Used only in the UK-facing versions. Cambridge's published
%          table for US degrees: 3.5 = 2:1, 3.7 = High 2:1, 3.9 = First.
%          3.60 is therefore a 2:1. Change this only if the GPA crosses 3.7.
%          To drop the label, delete "(\gpaUK)" where it is used.
\newcommand{\gpaNow}{3.61}
\newcommand{\gpaUK}{UK Good 2:1 Honours equivalent}
   % name, contact details \& GPA figures
% Tighter layout for this version only
\geometry{left=0.7in,right=0.7in,top=0.6in,bottom=0.75in}
\titleformat{\section}{\cvneed{4\baselineskip}\normalsize\scshape\bfseries}{}{0pt}{}[\vspace{-6pt}\textcolor{cvrule}{\rule{\linewidth}{0.5pt}}]
\titlespacing*{\section}{0pt}{7pt}{3pt}

\begin{document}

% ---------------------------------------------------------------------
% HEADER: no street address, photo or date of birth (UK norm).
% ---------------------------------------------------------------------
\cvheader{\cvname}{%
  \cvemail\sep\cvphone\sep\cvweb\sep ORCID \cvorcid}

% =====================================================================
\section{Education}
% Sources: NJIT transcript issued 10 Jun 2026; NJIT financial aid record
%   (2 Oct 2026). Official degree name (NJIT catalog): "B.S. in
%   Mathematical Sciences, Mathematical Biology Concentration".
%   The transcript prints Program \& Major as "Mathematical Sciences"
%   (Department: Mathematics) \& does not print the concentration.
% Minors: the transcript lists Biomedical Engineering Technology \&
%   Philosophy \& Applied Ethics. History is approved \& complete (History
%   dept, Nov 2025; confirmed after Spring 2026), but NJIT's system holds
%   only two minors, so the registrar adds the third at graduation. Until
%   then the transcript shows two.
% "Dean's List in all semesters": on the transcript for Fall 2023, Spring
%   2024, Fall 2024, Fall 2025 \& Spring 2026; the Lund term has no line.
% "over $100,000": the aid record shows $98,786 in merit scholarships
%   through Fall 2026. The Spring 2027 payment takes it past $100,000,
%   so this figure is right for anything submitted after that posts.
% Modules: every related module, alphabetical within each row, no
%   grades. \key{...} = bold. Lund rows use Lund's own course titles.
% Lund (official transcript, 6 Oct 2026): 44.0 ECTS. Second cycle:
%   Biomathematics (3), Immunotechnology (4), Microscopy, Bio-Imaging (G),
%   Ultrasound Physics \& Technology (4). First cycle: Patent \& IP
%   Rights (3), Swedish (B). Neuroengineering was not completed (labs \&
%   project part only, 3.5 of 7.5 ECTS), so it is not listed.
% ---------------------------------------------------------------------
\entry{B.S. Mathematical Sciences (Mathematical Biology)}{Sep 2023 -- May 2027}
      {Albert Dorman Honors College, New Jersey Institute of Technology (NJIT)}{Newark, NJ, USA}
\begin{cvitems}
  \item \gpaUK; Dean's List in all semesters
  \item Minors: Biomedical Engineering Technology; History; Philosophy \& Applied Ethics
  \item Merit scholarships (over \$100,000): Honors Merit Award; Honors Residential Scholarship; NJIT Academic Excellence Scholarship; Highl\&er Achievement Scholarship
\end{cvitems}
\begin{modules}
  \item[Biomedical] \key{Applications in Healthcare}; Biochemistry; Cell \& Molecular Biology
  \item[Medicine \& society] \key{Biomedical Ethics}; Health, Justice, \& Capitalism; Medical Sociology; \key{Medicine \& Law in Modern History}
  \item[Quantitative] Applied Numerical Methods; Differential Equations (ordinary \& partial); \key{Linear Algebra}; \key{Mathematical Modeling}; Stochastic \& Discrete Models in Biology
  \item[Independent study (elective individual modules with professor per their discretion)] Hawkes process model of amyloid nucleation (Dr James Maclaurin); \key{Simplyifying \& modeling the MAPK/ERK pathway for therapeutic intervention} (Dr Horacio Rotstein, in progress); American medicine \& end-of-life care (Dr Stephen Pemberton, in progress); Intellectual history in the early modern Khalsa (Harvard Sikh Center \& Satnam Singh, in progress)
\end{modules}

\entry{Exchange semester (Global E3), Lund University, Sweden}{Jan 2025 -- Jun 2025}
      {Faculty of Engineering (LTH); 44 ECTS}{}
\begin{modules}
  \item[Master's level] Biomathematics; Immunotechnology; Microscopy, Bio-Imaging; \key{Ultrasound Physics \& Technology}
  \item[Undergraduate] Patent \& Intellectual Property Rights
\end{modules}

% =====================================================================
\section{Research Experience}
% Ongoing research only. The self-contained independent studies (Maclaurin,
%   Rotstein, Pemberton) are in the Education modules block; the Harvard
%   Sikh Center fellowship is under Funding \& Honours.
% EMS study (emails, Nov 2025 - Apr 2026): "A Comprehensive Study of
%   Operational Models in University-Based EMS Organizations", PI Dr Chang
%   Yaramothu, run in BMET 311-312. Three students revised the IRB protocol,
%   Google Form questionnaire \& recruitment email; you sent the
%   recruitment request to NCEMSF's regional coordinators (13 Apr 2026).
%   Still grey: how many organisations responded, \& any finding. If the
%   study continues past May 2026, change the end date to Present.
% ---------------------------------------------------------------------
\entry{Undergraduate Researcher \& Computational Biology Team Lead}{Oct 2023 -- Present}
      {KumarLab, PI: Dr Vivek Kumar, University of Houston (previously NJIT)}{Houston, TX, USA}
\begin{cvitems}
  \item Design cyclic peptide therapeutics \& test via molecular dynamics simulation; my top-ranked design is the lead peptide of a published multiple sclerosis study [1]
  \item Led the male contraceptive project: wrote \& won funding for peptide to inhibit fertilization; manuscript \& patent application in preparation [2]
  \item In vivo work: conducted dissections \& helped facilitate animal work on angiogenic hydrogels that were later published
  \item Train \& lead a team of students (15+ trained); over \$15,000 in direct project funding
\end{cvitems}

\entry{Undergraduate Researcher (UROP), Computational Bionanotechnology}{Jun 2026 -- Present}
      {PI: Prof.\ Stefano Angioletti-Uberti, Department of Materials, Imperial College London}{London, UK}
\begin{cvitems}
  \item Use protein language models to design peptides that mimic natural lig\&s; contribute code to BAGEL, a tool for searching biomolecular energy l\&scapes; currently testing in-vitro
\end{cvitems}

\entry{Student Researcher, national study of collegiate EMS operating models}{Sep 2025 -- May 2026}
      {PI: Dr Chang Yaramothu, School of Applied Engineering \& Technology, NJIT}{Newark, NJ, USA}
\begin{cvitems}
  \item With two other students, wrote the IRB protocol, questionnaire \& recruitment materials for a national survey of the training \& operating capabilities of US collegiate EMS organisations; ran recruitment through National Collegiate EMS Foundation's regional coordinators; currently analysing data
\end{cvitems}

% =====================================================================
\section{Publications \& Presentations}
% [2]: your current CV writes the withheld target as "*****"; here the
% phrase is simply left out of the title.
% ---------------------------------------------------------------------
\cvsubsection{Publications}
\begin{publist}
  \item C. Heffernan, O. West, \me{K. Singh}, J. Dodd-o, M. James, A. Baradi, K. Rybak, A. Lateef, S. Elkabes, V. Kumar.
        ``A novel cyclic peptide that blocks human autoimmune antibodies for the treatment of multiple sclerosis.''
        \emph{Eur.\ J.\ Med.\ Chem.}\ (2026).
        \href{https://doi.org/10.1016/j.ejmech.2026.119173}{doi:10.1016/j.ejmech.2026.119173}
        \contrib{ran the molecular dynamics simulations; designed the published peptide}
  \item \me{K. Singh}, C. Heffernan, O. West, J. Dodd-o, V. Kumar.
        ``Machine learning for rational design of a cyclic peptide-based male contraceptive.''
        \status{In preparation}; patent application in preparation
  \item J. Lau, D. Yun, \me{K. Singh}, A. Patananan, N. Jadavji.
        ``Academic publishing as a launchpad: how publishing research in an undergraduate journal shapes future careers.''
        \status{In preparation}
\end{publist}

\cvsubsection{Selected presentation}
\pres{\me{K. Singh}, C. Heffernan, J. Dodd-O, V. Kumar}{Development of a novel male contraceptive using cyclic peptides to inhibit fertilization}
     {Posters: AAP/ASCI/APSA Joint Meeting (2026); National Contraception Meeting (2025); Biomedical Engineering Society (2024); NJIT Honors Interdisciplinary Research Forum (2025, \textbf{First Place})}
% Left off to hold two pages (both are on the general CV): the CB1
% receptor poster (UH Undergraduate Research Day, 2026) \& the journal
% workflow poster (NCUR, 2026).

% =====================================================================
\section{Research Funding \& Honours}
% Grants won for a project first, then selective programmes \& prizes.
% Amounts: only the Grace Hopper figure is on the aid record; add the
% others if you want numbers here.
% ---------------------------------------------------------------------
\oneline{\textbf{UROP research bursary}, Imperial College London -- protein language models}{2026}
\oneline{\textbf{NSF I-Corps Regional Program}, Northeast Hub -- customer discovery for peptide therapeutics}{2026}
\oneline{Emerging Scholars Research Fellowship, Harvard Sikh Center}{2026}
\oneline{\textbf{Male Contraceptive Initiative} project grant \& travel award -- cyclic peptide contraceptive}{2025}
\oneline{Moonshot Grant for Engaged Scholarship, NJIT Honors College -- \textbf{one recipient per annum}}{2025}
\oneline{\textbf{Grace Hopper AI Research Institute} award, NJIT -- protein design course; \$3,330}{2025}
\oneline{{Research seed grant \& Honors Summer Research Institute} grant, NJIT}{2024}

% =====================================================================
\section{Clinical \& Translational Experience}
% The bedside \& industry ends of "bench to bedside".
% The startup is left unnamed here because you had described it as
% stealth. Your site \& current CV do name it; put the name back if
% that is fine.
% ---------------------------------------------------------------------
\cvsubsection{Clinical \& public health}
\entry{Emergency Medical Technician (New Jersey licence)}{Apr 2022 -- Present}
      {NJIT First Aid Squad (Public Relations Chair); Marlboro First Aid Squad (Cadet Crew Chief)}{}
\begin{cvitems}
  \item over 2000 hours h\&ling emergency medical situations as a licensed medical professional
\end{cvitems}
\oneline{\textbf{Stop the Bleed Instructor}, NJIT \& New Jersey gurdwaras (English--Punjabi)}{2025 -- Present}
\begin{cvitems}
  \item taught first aid skills to 100+ as certified instructor, supported through the first aid squad \& the NJ Office of Homel\&
Security \& Preparedness (NJOHSP) respectively
\end{cvitems}

\cvsubsection{Industry \& commercialisation}
\oneline{\textbf{Co-founder \& technical lead}, Axioma Therapeutics -- funded through NSF I-Corps program}{2026}
% Left off to hold two pages: NSF ARTISAN membership (on the general CV).

% =====================================================================
\section{Leadership, Teaching \& Mentoring}
% Chosen by scale \& by who benefited.
% ---------------------------------------------------------------------
\oneline{\textbf{22\textsuperscript{nd} Editor-in-Chief}, Journal of Young Investigators {oldest international undergraduate journal; increase access to undergraduate research communication [3]; funded \& featured by Nature, Science, NSF, Burrough's Wellcome Fund, etc.}}{2024 -- Present}
\oneline{\textbf{Research Chair}, American Association of Physicians of Indian Origin PreMed Board}{2025 -- Present}
\oneline{\textbf{Research Lead}, Governor's STEM Scholars -- mentor a high school research team}{2026 -- 2027}
\oneline{\textbf{Teaching}: research ethics, Science Mentorship Institute; Sikh history, Khalsa School}{2020 -- Present}

% =====================================================================
\section{Research Skills \& Training}
% Grey = not stated anywhere I could check. Name tools, not adjectives.
% ---------------------------------------------------------------------
\begin{skills}
  \item[Computational:] Rosetta; PyMOL; GROMACS; NAMD; VMD; OpenMM; protein language models; AlphaFold; Boltz; ESMFold; CHARMM
  \item[Laboratory:] SPR, general dissection (amateur)
  \item[Data analysis \& statistics:] Python; RosettaScripts; GROMACS Scripting; \LaTeX; RMSD
  \item[Clinical research training:] EMT licence; AHA BLS \& CPR; CITI Biomedical Research \& RCR
  \item[Languages:] English (C2); Punjabi (C1); Swedish (B2); Spanish (B1); Russian (B1)
\end{skills}

\end{document}